\section{Results}
\label{sec:results}

\paragraph{Where does the dataset effect appear?}
If the datasets differed only in noise, the retention gap would scatter around zero across the grid.
Instead, bit width and group size set the level of retention, and the gap concentrates where retention lies between its two limits (\cref{fig:gap}).
Averaged over the nine backbone and regime settings, retention is at most \V{ret_w1_max} at 1 bit and at least \V{ret_w8_min} at 8 bits for every group size.
At W3/G128, the \SC{} group retains \V{retSc_w3g128} against \V{retOther_w3g128} for the other group, a gap of \V{gap_w3g128}.
The gap is positive in \V{gapPos_w3} of the \V{nSettingGs} combinations of setting and group size at 3 bits, and in \V{gapPos_w4} of \V{nSettingGs} at 4 bits.
Smaller groups move this band to fewer bits: the gap is \V{gap_w2g8} at W2/G8 but only \V{gap_w2g128} at W2/G128, where mean retention has fallen to \V{ret_w2g128}.
\textbf{The dataset effect is therefore a property of the transition band, located by bit width and group size together.}
At 1 bit the gap is negative for BEiTv2 FT, also after chance correction; several of its \SC{} baselines are near chance, where both measures are unstable.
The gap at W3/G128 survives chance-corrected retention (\V{gapChance_w3g128}), removal of the \V{lowBaseline} baselines below 15\% (\V{gapNoLow_w3g128}), and removal of any one \SC{} dataset (at least \V{looMinGap}), staying positive in all nine settings (\cref{tab:robust}).
A bootstrap over datasets places its 95\% interval above zero in \V{bootPositive} of nine settings, all but BEiTv2 FT.


\paragraph{Is the dataset ranking stable?}
A dataset property that drives difficulty should rank the datasets similarly across conditions and models.
Within each setting the four primary conditions agree well, with Kendall's $W$ between \V{wMin} and \V{wMax} (\cref{tab:gates}).
Across settings the agreement is weaker: the median Spearman correlation between two settings' scores is \V{agreeMedian}, ranging from \V{agreeMin} to \V{agreeMax} (\cref{fig:agreement}).
BEiTv2 CAL and BEiTv2 FT agree least, so the backbone alone does not fix the ranking.
Still, the \SC{} datasets hold at least \V{scHardestMin} of the five hardest places in every setting.

\begin{table}[t]
\centering
\caption{Difficulty score of the {\V{gridDatasets}} complete datasets. $W$: agreement of the four primary conditions. $\rho$: Spearman correlation of the score with the baseline margin $q_{10}$ (one-sided permutation $p$) and with FP top-1. Err: leave-one-dataset-out squared error ($\times 10^{3}$) of the controls (Ctrl.) alone and with $q_{10}$ added.}
\label{tab:gates}
\setlength{\tabcolsep}{1.8pt}
\small
\begin{tabular}{llrrrrrr}
\toprule
 & & & \multicolumn{2}{c}{Margin $q_{10}$} & FP top-1 & \multicolumn{2}{c}{Err $\downarrow$} \\
\cmidrule(lr){4-5}\cmidrule(lr){7-8}
Backbone & Regime & $W$ $\uparrow$ & $\rho$ & $p$ & $\rho$ & Ctrl. & +$q_{10}$ \\
\midrule
ViT-B/16 & FT & 0.89 & $-$0.98 & $<$0.01 & $-$0.97 & 16.3 & 62.5 \\
ViT-B/16 & FZ & 0.93 & $-$0.92 & $<$0.01 & $-$0.90 & 15.9 & 184.8 \\
ViT-B/16 & CAL & 0.80 & $-$0.83 & $<$0.01 & $-$0.83 & 21.8 & 21.7 \\
Swin-B & FT & 0.81 & $-$0.86 & $<$0.01 & $-$0.78 & 37.2 & 74.5 \\
Swin-B & FZ & 0.96 & $-$0.97 & $<$0.01 & $-$0.98 & 11.9 & 164.4 \\
Swin-B & CAL & 0.76 & $-$0.63 & $<$0.01 & $-$0.59 & 73.8 & 84.7 \\
BEiTv2-B & FT & 0.83 & $-$0.48 & 0.04 & $-$0.40 & 120.1 & 111.8 \\
BEiTv2-B & FZ & 0.93 & $-$0.92 & $<$0.01 & $-$0.95 & 11.1 & 86.1 \\
BEiTv2-B & CAL & 0.94 & $-$0.83 & $<$0.01 & $-$0.74 & 49.5 & 32.5 \\
\bottomrule
\end{tabular}
\end{table}


\paragraph{Do baseline margins explain the difficulty?}
If quantization perturbs logits by a similar amount everywhere, datasets with many small-margin images should lose more.
The prespecified margin signal holds in \V{gatePass} of nine settings, with a median $\rho$ of \V{rhoMarginMedian} and $p<0.05$ in all \V{gatePPass} (\cref{tab:gates}).
BEiTv2 FT fails only the strength criterion, with a $\rho$ of \V{rhoMarginBeitFt}.
However, baseline accuracy correlates with the score almost as strongly, with a median $\rho$ of \V{rhoAccMedian}.
Adding $q_{10}$ to the controls lowers the held-out error in only \V{marginHelps} of nine settings.
Controlling for accuracy, the partial correlation of $q_{10}$ with the score stays negative in \V{partialNeg} of nine settings, with a median of \V{partialMedian}.

Baseline accuracy also cannot fix the retention of an individual context (\cref{fig:accuracy}); over the \V{pooledN} contexts at W3/G128 it explains \V{pooledR2}\% of the variance in a linear fit.
On BEiTv2 CAL, Cars and Food start from \V{pairALowFp}\% and \V{pairAHighFp}\% but retain \V{pairALowR} and \V{pairAHighR}.
On Swin-B CAL, SoyGene and iNat17 start from \V{pairCLowFp}\% and \V{pairCHighFp}\%, and the soybean dataset retains \V{pairCLowR} against \V{pairCHighR}.
In both pairs the same dataset retains less in all four primary conditions.
\textbf{Baseline accuracy explains much of the difficulty, but not all of it.}

\begin{figure}[t]
  \centering
  \includegraphics[width=\linewidth]{figures/accuracy_retention.pdf}
  \caption{Retention at W3/G128 against baseline top-1 for the \V{pooledN} contexts of the \V{gridDatasets} complete datasets. Lines join three pairs from one backbone and regime with nearly equal baseline accuracy.}
  \label{fig:accuracy}
\end{figure}

\begin{table}[t]
\centering
\caption{Retention gap (other $-$ soy+cotton) at W3/G128 before and after matching the baseline margin distributions, bin width 0.05. Red.: reduction of the gap. Overlap is the shared histogram mass; ESS is the lower of the two groups' matched-to-raw effective sample size.}
\label{tab:margin}
\setlength{\tabcolsep}{1.6pt}
\small
\begin{tabular}{llrrrrr}
\toprule
Backbone & Regime & Raw & Matched & Red. (\%) & Overlap & ESS (\%) \\
\midrule
ViT-B/16 & FT & 0.212 & 0.042 & 80.0 & 0.26 & 23.4 \\
ViT-B/16 & FZ & 0.310 & 0.188 & 39.2 & 0.40 & 42.6 \\
ViT-B/16 & CAL & 0.268 & 0.118 & 55.9 & 0.24 & 29.6 \\
Swin-B & FT & 0.498 & 0.248 & 50.2 & 0.23 & 26.3 \\
Swin-B & FZ & 0.569 & 0.412 & 27.6 & 0.34 & 32.9 \\
Swin-B & CAL & 0.341 & 0.112 & 67.1 & 0.26 & 31.9 \\
BEiTv2-B & FT & 0.408 & 0.188 & 53.9 & 0.19 & 20.9 \\
BEiTv2-B & FZ & 0.184 & 0.078 & 57.4 & 0.37 & 34.1 \\
BEiTv2-B & CAL & 0.325 & 0.061 & 81.3 & 0.34 & 47.5 \\
\bottomrule
\end{tabular}
\end{table}


Margin matching asks the same question within the W3/G128 test images, where the \SC{} images concentrate at small margins (\cref{fig:margin}).
Matching both groups to their common margin distribution reduces the gap by \V{reductionMin}\% to \V{reductionMax}\% (\cref{tab:margin}).
The matched gap stays positive in \V{matchedPositive} of nine settings, so the \SC{} images lose more even at equal margin.
The comparison covers only the shared \V{overlapMin} to \V{overlapMax} of the histogram mass, and the matched ESS falls to as little as \V{essRatioMin}\% of the raw ESS.
Matching aligns only the distance to the second-ranked class, yet a median of \V{beyondSc}\% of the \SC{} swaps, against \V{beyondOther}\% for the other datasets, end on a class the baseline ranked below second.
This share follows the margin ($\rho$ \V{beyondRhoMargin} with $q_{10}$) rather than the class count ($\rho$ \V{beyondRhoClasses}): Cotton with \V{cottonClasses} classes reaches \V{beyondCotton}\%, against \V{beyondCub}\% for CUB with \V{cubClasses}.
The top-1 to top-2 margin therefore misses most \SC{} swaps, consistent with the positive matched gap, and margins to lower-ranked classes are the natural next diagnostic.

\paragraph{Do feature and decision changes explain it?}
Quantized features of the hard datasets do not drift further from the baseline.
At mid depth, the median CKA between baseline and W3/G128 features is \V{midCkaSc} for the \SC{} group and \V{midCkaOther} for the other group (\cref{fig:cka}), and final-layer CKA correlates with the score at a median $\rho$ of only \V{rhoCkaMedian} (\cref{tab:properties}).
Feature similarity flags one collapsed setting, BEiTv2 FZ, whose median final CKA is \V{beitFzCka} and whose mean accuracy falls from \V{beitFzFp}\% to \V{beitFzQ}\% (\cref{tab:w3}).
At the decision level, damage exceeds rescue by a factor of at least \V{damageOverRescueMin} in every setting.
Cross-fitted class-offset correction recovers \V{biasGainMin} to \V{biasGainMax} points of top-1, while shuffled-class controls lose \V{shuffleLossMin} to \V{shuffleLossMax} points, so part of the damage is a class-specific logit shift.
Counts of competing classes improve damage prediction beyond the margin in \V{competitorBetter} of nine settings.
